\chapter{Grounding the Tower, and the First Record}
\label{chap:grounding-and-first-record}

The construction now has a concrete operator and a real norm; it can stop reasoning about apparatus in
the abstract and ground the whole tower of characteristics from Part I on the actual three-node form it
built two chapters ago. Doing so brings the book's longest-running discipline --- the careful
separation of what is demonstrated from what is merely named --- to a head, and turns it from a habit of
exposition into a theorem. This chapter does three things. It grounds the abstract apparatus on the
concrete operator. It proves that whether a quantity counts as physical or as metaphysical is not
absolute but relative to the apparatus asked to measure it. And it writes down the construction's first
genuine measurement record, a single invariant that is physical to one apparatus and metaphysical to
another, stated entirely in the apparatus's own vocabulary.

\section{The funge and the crank}
\label{se:the-funge-and-the-crank}

Everything in this chapter rests on a distinction the construction has used informally since its first
page and now makes exact. When the apparatus reads a Fact, the reading goes one of two ways. Read
forward, with the truth taken as it stands, the reading is a \emph{funge} --- and counting the funges is
the measurement, the tally of matter the apparatus actually registers. Read reversed, with the truth
turned against itself, the reading is a \emph{tange}, and a tange does not add to the count; it turns
the crank. The image is exact and a little unkind: an apparatus that produces only tanges is a crank in
the colloquial sense too, a device that turns and turns and registers nothing, its motion all residue
and no matter. A genuine measurement is a count of funges; the tanges are the leftover, the part of the
reading that spins the mechanism without recording anything. The balance between the two --- how many
readings come out forward and how many reversed --- is a parity, and parity is the quantity this whole
chapter turns on.

It helps to make the parity concrete, because it is the hinge of everything that follows. Lay out the
apparatus's readings as a tape of bits, one per Fact: a one where the reading came out forward, a funge,
and a zero where it came out reversed, a tange. The measurement is the count of ones --- the funges, the
matter --- and the count of zeros is how many times the crank turned. Two tapes that count the same
number of funges register the same amount of matter, however their tanges are arranged; the matter is
the forward count and nothing else. Parity is the simplest fact about such a tape: whether the count of
reversed readings is even or odd. And parity is conserved in a way the individual readings are not ---
one can rearrange which Facts come out forward and which reversed, but certain rearrangements are
forbidden because they would change the parity, and a change of parity is one the apparatus cannot make
without an external source for the missing turn. That conservation is what turns parity from a
bookkeeping curiosity into an obstruction with teeth.

This recasts the book's central activity in two words. To measure is to funge: to read truths forward
and count them, pooling the matter the apparatus can register into a number. The residue, the part that
will not come out forward, turns the crank and is not counted. The construction has, in effect, been
distinguishing funges from cranks since it first asked how true a single Fact is, but only now, with a
concrete operator under it, can it say precisely which readings are matter and which are merely motion.
The measurement is the funge count; everything else is the crank turning.

\section{Physical and metaphysical}
\label{se:physical-and-metaphysical}

With the funge made precise, the construction can finally say what it means for a quantity to be
physical. An invariant is \emph{physical} relative to a class of apparatus when some apparatus in the
class realizes it --- when there is a device that counts it out in funges, that registers the invariant
as matter rather than spinning past it. It is \emph{metaphysical} relative to that class when no
apparatus in the class can do so, when every device that might measure it is blocked, its motion all
crank and no count. And the construction proves the blocking can be real and certified: there are
invariants for which every apparatus in a given class provably fails, the failure sourced from a genuine
obstruction --- a parity that cannot be balanced, an orientation that cannot be closed. An odd invariant,
in the simplest case, has no witness in a class whose apparatus can only register even counts; it is
metaphysical there, not because it is meaningless but because nothing in that class can read it.

The relativity is easiest to see in the parity case the construction proves. Suppose an apparatus can
only ever register readings in matched pairs --- its construction forces an even number of forward
reads, so the parity it can produce is always even. Hand it an invariant whose honest measurement
requires an odd count, a single unpaired funge, and the apparatus cannot deliver it: every configuration
it can form has even parity, so the odd invariant has no witness anywhere in its repertoire. The
invariant is not false and not meaningless; it is simply unreadable by this apparatus, metaphysical
relative to it. Now enrich the apparatus --- give it one more degree of freedom, a single reading that
can supply the missing odd turn --- and the same invariant is suddenly realizable: there is now a
configuration that counts it out, an unpaired funge the richer apparatus can form. The invariant did not
change. The apparatus did, and with it the invariant's status, from metaphysical to physical, by the
addition of exactly the parity that had been missing.

The decisive theorem is that this status is relative. The very same invariant that is metaphysical to
one class of apparatus is realized --- counted out in funges, made physical --- by a richer class that
has the parity the poorer one lacked. Physical and metaphysical are not properties an invariant carries
in itself; they are relations between an invariant and an apparatus, and an invariant moves from one to
the other as the apparatus is enriched. This is the experimental library's own honesty made into
mathematics. The experiment under Wittgenstein's name draws the line where the construction draws it:
what can be said is what the apparatus can register, and what lies past its registering is not false but
unsayable in its language --- metaphysical relative to it, awaiting a richer apparatus in which it can
be spoken. The experiment on the topological integer count supplies the obstruction's shape: an
orientation, a parity, an integer that an apparatus either can or cannot close, and that decides whether
the invariant is matter to it or silence. The book's long discipline of marking what is demonstrated
against what is assumed is, it turns out, exactly this relation. A quantity demonstrated is a quantity
some apparatus funges; a quantity assumed is one metaphysical relative to the apparatus that names it,
carried as a label because no device in reach can count it out.

This is worth stopping on, because it changes the status of everything the book has done. Every chapter
has ended by separating what it demonstrated from what it assumed, and that separation has looked like a
matter of intellectual honesty --- a discipline imposed from outside, a promise to the reader. This
chapter shows it is not a promise but a theorem. A demonstrated quantity is one the construction's
apparatus funges; an assumed one is metaphysical relative to that apparatus, named because it cannot be
counted out. The line between the two columns of every chapter's closing note was never the author's
scruple; it was the parity obstruction, deciding case by case which invariants the apparatus could read
and which it could only point at. The book's honesty, in other words, is forced by the same mathematics
the book is about, and a register-honest closing note is simply a measurement of which side of the
obstruction each claim falls on.

\section{The first record}
\label{se:the-first-record}

Now the construction writes its first genuine measurement record, and it is fitting that the first thing
it records is an instance of exactly the relativity just proved. Take the rigid mode of the three-node
operator --- the uniform shift that the unanchored form cannot see. Relative to the unanchored apparatus
it is physical: the unanchored energy has a null mode, a nonzero configuration the energy reads as zero,
and that null mode is a real feature the apparatus registers. Relative to the anchored apparatus it is
metaphysical: the anchored energy, definite, has no null mode at all, so the same rigid configuration
simply does not exist as a thing the anchored apparatus can read. One invariant, physical to the
apparatus without the anchor and metaphysical to the apparatus with it --- the abstract relativity of
the previous section, instantiated in the smallest concrete operator the construction owns.

The worked case is the one the construction commits to its ledger first. The rigid mode is the
configuration in which all three nodes are displaced together, the uniform shift. To the unanchored
apparatus this is a genuine object: the unanchored energy, blind to uniform shifts, reads it as a null
mode --- a nonzero configuration of zero energy --- and a null mode is something the apparatus registers,
a feature it can point to and count. To the anchored apparatus the same uniform shift is nothing at all:
the anchor forbids it, the anchored energy has no null mode, and there is simply no configuration there
for the rigid mode to be. So the question \emph{is the rigid mode real?} has no apparatus-independent
answer. It is real to the apparatus without the anchor and absent to the apparatus with it, and the only
honest report is the two-sided one --- which is exactly what the construction records.

That two-sided fact is the first record. The construction proves it as a single conjoined statement ---
the anchored energy has no null mode, and the unanchored energy does --- and, crucially, it states the
record in the apparatus's own vocabulary, in terms of null modes and energies and anchors rather than in
any outside language imported to describe them. This is what it means for the construction to have
become self-describing in earnest: it can write down a measurement whose entire content is a relation
between an invariant and its apparatus, using only the words the apparatus itself supplies. The first
record is small --- it concerns a single rigid mode on three nodes --- but it is the template for every
measurement the rest of the book makes. To record a measurement is to state, in the apparatus's own
terms, which invariants it funges and which it cannot, and the rigid mode of the three-node form is the
first thing the construction commits to that ledger.

That the record is stated in the apparatus's own vocabulary is not a stylistic nicety but the whole
achievement. A measurement reported in some external language --- in the experimenter's intentions, in a
metaphysics of what is really there --- would smuggle in exactly the apparatus-independent standpoint the
chapter has just shown does not exist. By writing the first record using only null modes, energies, and
anchors --- the apparatus's own terms --- the construction makes a measurement that needs no outside
vantage to be understood and no outside authority to be trusted. It is a record a different reader,
holding the same apparatus, would write identically, because it says nothing the apparatus does not
itself supply. Self-description, reached in principle at the first theorem of the eighth chapter, here
becomes self-measurement: the construction not only proves things about its own apparatus but records
readings taken with it, in its words, answerable to nothing but it.

\section{What is demonstrated, and what is assumed}
\label{se:demonstrated-assumed-ch17}

This chapter's accounting is, unusually, partly about accounting itself.

What is demonstrated is the grounding, the relativity, and the record. The abstract tower of Part I is
grounded on the concrete three-node operator, the self-bootstrap of the eighth chapter realized on an
actual form. The physical-versus-metaphysical distinction is proved relative to the apparatus class:
some invariants are metaphysical relative to a poorer class, sourced from a genuine parity obstruction,
and the same invariants are realized by a richer one. And the first record is proved and stated
internally: the rigid mode has no null mode under the anchor and a null mode without it. Each is a
theorem about finite objects.

\begin{quote}
\emph{A note on the labels --- and on this book's own method.} The reading offered here is unusually
self-referential, and worth marking as such. To call the forward reading \emph{matter} and the reversed
one \emph{crank-turning}, and to align \emph{physical} with what an apparatus funges and
\emph{metaphysical} with what it cannot, is to give the book's running distinction between demonstrated
and assumed a precise relational form. That alignment is itself a modeling choice --- the mathematics is
the parity obstruction and the relativity theorem; the identification of ``physical'' with
``apparatus-realizable'' is the interpretation laid over it. But it is the interpretation the whole book
has been operating under, now stated openly: nothing in this construction is physical absolutely;
things are physical to apparatus, and the honest report always names the apparatus.
\end{quote}

What is assumed remains the standing oracle and the modeling identification just marked --- nothing new
of substance. With the tower grounded, the physical and the metaphysical sorted by apparatus, and the
first record written in the construction's own words, the apparatus side of the account is essentially
finished. The next chapter sets the apparatus in motion as an integrator, producing not just a verdict
but a certificate of its measurement, and from there the construction turns to read, off this machinery,
the first object it dares to give a particle's name.
